\chapter{Glossary}\label{book-glossary}

Terms in this glossary describe the teaching companion and the vocabulary used to explain it. An identifier in a historical record can have a different version-specific meaning.

\subsection{Adapter}

Code that connects a host or framework boundary to application contracts. The companion has keyboard, scheduler and shutdown adaptations.

\subsection{App}

The Bijou runtime contract implemented by init, update and view in the companion. AppModel includes viewport state around the task model.

\subsection{Application model}

A snapshot of facts used to decide behavior and presentation. The companion separates the Monitor task model from viewport dimensions and scroll offset.

\subsection{Cancellation}

A request to stop owned work. It manages resources; it does not by itself establish that a queued result can no longer arrive.

\subsection{Cell}

A unit in a terminal Surface. A code point, grapheme cluster and occupied terminal width are different concepts.

\subsection{Checkpoint}

A reproducible companion command that exposes a selected behavior of the finished implementation. It is not a separate historical code release.

\subsection{Cmd}

A framework command used by the App to perform work outside a pure transition or request runtime behavior such as quit.

\subsection{Effect}

An application description of work to request, such as starting a particular run or loading logs for a view epoch.

\subsection{Epoch}

A generation number for the current view request. A completion must match the selected job, screen and epoch before becoming visible detail.

\subsection{Focus and input ownership}

The policy determining which active surface receives an intention. Drawing a dialog does not establish this policy.

\subsection{Grapheme cluster}

A user-perceived text unit that can contain multiple code points. The companion wraps clusters and measures their terminal width.

\subsection{Historical record}

A preserved source document whose status and assertions retain their original context. It is not a current implementation guarantee.

\subsection{Identity}

A stable value used to address a job independently of its current position or display label.

\subsection{Immutable snapshot}

A model value that does not change after construction. Later transitions produce another snapshot.

\subsection{Intention}

An application event expressing a requested action, such as start, move or dismiss, after interpreting framework input.

\subsection{KeyMsg}

A decoded framework keyboard message translated by the companion App into an application intention.

\subsection{Mode}

A presentation and delivery choice. The companion offers interactive, static, pipe and accessible modes with distinct output contracts.

\subsection{Port}

An explicit interface to an external capability. The scheduler port permits deterministic tests and a host timer implementation.

\subsection{Run identity}

A generation identifying one attempt at a job. It prevents a completion from an earlier attempt updating a restarted job.

\subsection{Semantic styling}

Choosing appearance from a content role while keeping application meaning independent of color or other decoration.

\subsection{Stale result}

A result whose original owner or generation no longer matches the current model. The companion rejects stale job and view completions.

\subsection{Surface}

Bijou representation of rendered terminal cells used by the companion view and controlled rendering observations.

\subsection{TEA}

The Elm Architecture vocabulary for the state, update and view loop with explicit work outside the transition. The book uses the pinned App contract to make that vocabulary concrete.

\subsection{Transition}

The result of applying an event to a model: a next model and any effect descriptions.

\subsection{Viewport}

The width and height available for the current view, together with the application decisions governing what is visible.
